\subsection{Mappings}
\label{sec:mappings}

Let $L=\min(R,C)$. All mappings below are those of the library; each is
loaded with one \cmd{CFG} per call (per column tile for matrix kernels).

\emph{Element-wise (diagonal).} Vector $a$ enters from the north and $b$
from the west; PEs above the diagonal pass north values south, those below
pass west values east, and the $L$ diagonal PEs apply the operation.
Because operands need up to $L-1$ hops, one chunk takes $L$ steps and yields
$L$ results. Inputs stay fixed during the chunk, so the diagonal PEs
recompute the same value until the last step. Partial chunks are padded with
zeros, and only valid outputs are returned.

\emph{Dot product.} PE$(0,0)$ executes $r\leftarrow r+a_ib_i$ on its edge
inputs; the host clears it once and supplies one pair per step. The partial
sum stays local: local output-stationary accumulation is supported, and the
other $RC-1$ PEs are idle.

\emph{Inclusive scan.} PE$(0,c)$ computes $\mathrm{op}(\mathrm{west},
\mathrm{north}_c)$ and passes the fold east; the previous block's last
result enters west input~0 as a carry (the identity for the first block).
One block of $C$ elements takes $C$ steps: PE$(0,c)$ settles at step $c+1$
and recomputes thereafter. ADD, MUL, MIN and MAX folds are supported.

\emph{Matrix--vector ($y=Ax$).} For each column tile, row~0 holds $x_c$ in
its immediates and multiplies the matrix row injected from the north;
rows $1,\dots,R-1$ pass products south (Fig.~\ref{fig:systolic}). After $R$
steps the array holds $R\times C$ products, which the host reads, sums
across columns and accumulates across column tiles. On a full tile, $C$ PEs
multiply and $(R-1)C$ act as delay cells; the host reduction is part of the
algorithm. Column tiles whose $x$ entries are all zero, and all-zero
matrix tiles, are skipped. Convolution uses the same engine on implicitly
generated Toeplitz tiles~\cite{golub-vanloan}, visiting only tiles that
intersect the band.

\begin{figure}[t]
  \centering
  % One column of a 4x4 tile: products enter at row 0, move south one row per
% step, and are read together after R steps. Cells are drawn on one grid, so
% every edge is a single line.
\begin{tikzpicture}[x=10mm, y=7mm, font=\scriptsize]
  % cell (column t, row r): value; products shaded, zeros plain
  \foreach \t/\r/\v in {1/0/p_0, 2/0/p_1, 3/0/p_2, 4/0/p_3,
                        2/1/p_0, 3/1/p_1, 4/1/p_2,
                        3/2/p_0, 4/2/p_1,
                        4/3/p_0} {
    \fill[cReg] (\t-1,-\r) rectangle (\t,-\r-1);
    \node at (\t-0.5,-\r-0.5) {$\v$};
  }
  \foreach \t/\r in {1/1, 1/2, 2/2, 1/3, 2/3, 3/3}
    \node[text=cWire!70] at (\t-0.5,-\r-0.5) {$0$};
  % grid in figure units (TikZ's grid step would be in cm)
  \foreach \x in {0,...,4} \draw[cWire!45] (\x,0) -- (\x,-4);
  \foreach \y in {0,...,4} \draw[cWire!45] (0,-\y) -- (4,-\y);
  \draw[cRegE, very thick] (3,0) rectangle (4,-4);

  \foreach \t in {1,...,4} \node[lbl, above] at (\t-0.5,0) {$t=\t$};
  \foreach \r/\role in {0/{row 0: MUL}, 1/{row 1: PASS}, 2/{row 2: PASS}, 3/{row 3: PASS}}
    \node[lbl, anchor=east] at (-0.15,-\r-0.5) {\role};
  \draw[wire, cRegE] (4.25,-0.5) -- node[lbl, right, align=left, text=cRegE]
    {products\\move south} (4.25,-3.5);
  \node[lbl, below, align=center] at (2,-4.1)
    {initial registers zero;\quad last column: one read-back snapshot};
\end{tikzpicture}

  \caption{One column of a $4\times4$ matrix tile. Here $p_i=A_{i,c}x_c$.
  Time increases from left to right; the last column is the read-back snapshot.
  The host sums snapshot PE rows across all columns to obtain output rows in
  reverse order.}
  \label{fig:systolic}
\end{figure}


\emph{Forwarding limit.} An interior PE cannot keep a sum $s_t$ and forward
two operands $a_t$ east and $b_t$ south in the same step: after a MAC its
only visible register holds $s_{t+1}$. A one-MAC-per-PE-per-step
output-stationary schedule therefore needs extra registers, routing PEs,
time multiplexing or reconfiguration. This is a limit of that schedule on
this PE, not an impossibility result for output-stationary
mappings~\cite{sze-survey,eyeriss}.
